%%%%%%%%%%%%%%%%%%%%%%%%%%%%%%%%%%%%%%%%%%%%%%%%%%%%%%%%%%%%%%%%%%%%%%%%%%%%%%%%
%2345678901234567890123456789012345678901234567890123456789012345678901234567890
%        1         2         3         4         5         6         7         8

\documentclass[letterpaper, 10 pt, conference]{ieeeconf}

\IEEEoverridecommandlockouts
\overrideIEEEmargins

\input{usepackage}

\title{\LARGE \bf
Week 4 Initial Results: Data-Driven Priors and Physics-Based
Corrections \\ A Hybrid Flow Matching Approach to Safe Robot Motion
Planning}

\author{ \parbox{3 in}{\centering Zihao Sun}
        \hspace*{ 0.1 in}
        \parbox{2.5 in}{ \centering Charles Beck}
\thanks{UCLA ECE M237 Spring 2026 Course Project Milestone Reports.
Week 4---Initial Results.}
}

\begin{document}

\maketitle

\section{Main Idea}

We investigate how to merge an \textbf{offline}-trained generative
trajectory prior (learned from demonstration data) with an
\textbf{online} model-based score correction (computed analytically
from the known dynamics and cost) and study whether the resulting
hybrid yields trajectories that are simultaneously diverse and
physically feasible.

\section{Related Work}
\subsection{Overview}

Our project sits at the intersection of \emph{model-based score
computation} and \emph{generative trajectory planning}, with a
secondary connection to \emph{safe generative control}.
MBD~\cite{mbd} derives the score analytically from known dynamics
and cost, guaranteeing feasibility without any learned component;
its authors note data ``can be naturally integrated'' but leave
this unformalized---the gap we address. Diffusion-based
planners~\cite{janner2022} and flow-matching
methods~\cite{flowmp} learn expressive multimodal priors from
demonstrations but provide no feasibility guarantees out of
distribution. For collision avoidance we plan to incorporate
CBF-based score terms~\cite{safeflow}.

\subsection{Highlighted Literature Review}

Pan et al.~\cite{mbd} propose MBD, which replaces the learned
score of standard diffusion planners with one computed analytically
from $x_{k+1}=f(x_k,u_k)$ and $J(\tau)$ via Tweedie's formula and
self-normalized importance sampling. MBD defines a Boltzmann
trajectory distribution $p^{\mathrm{mod}}(\tau)\propto
\exp(-J(\tau)/\lambda)$, noises it via a VP-SDE, and approximates
the posterior mean by cost-reweighted rollouts; on Maze2D, Block
Push, and Franka manipulation benchmarks it matches or outperforms
Diffuser~\cite{janner2022} while requiring zero demonstration data.
Its main limitations are that (i) it cannot capture multimodal,
demonstration-informed behaviors, (ii) it assumes an accurate
dynamics model, and (iii) its importance-sampling estimator
suffers high variance when $\lambda$ is small or the horizon is
long. Our work targets all three by anchoring the denoising process
to a learned data prior and letting an adaptive $\lambda(x,t)$
govern the trust between the two components.

\section{Problem Formulation}

Let $\tau=(x_0,u_0,\ldots,x_{H-1},u_{H-1},x_H)$ be a trajectory of
the discrete-time system $x_{k+1}=f(x_k,u_k)$ with running cost
$J(\tau)=\sum_{k=0}^{H-1}\ell(x_k,u_k)+\ell_H(x_H)$. Given a
demonstration dataset
$\mathcal{D}=\{\tau^{(i)}\}_{i=1}^{N}\!\sim\!p^{\mathrm{data}}$,
our goal is to sample from the product-of-experts target
$p^{\mathrm{target}}(\tau)\propto p^{\mathrm{data}}(\tau)\,
[p^{\mathrm{model}}(\tau)]^{\lambda}$, where
$p^{\mathrm{model}}(\tau)\propto\exp(-J(\tau)/\beta)$ is the
Boltzmann distribution induced by the known dynamics and cost. At
each diffusion time $t$ this is equivalent to denoising with the
composite score
\begin{equation}
  \nabla\log p^{\mathrm{target}}_t \approx
  \nabla\log p^{\mathrm{data}}_t
  + \lambda(x,t)\,\nabla\log p^{\mathrm{model}}_t,
  \label{eq:score}
\end{equation}
which is the central object our project formalizes and analyzes.

\section{Sketch of Method}
\subsection{Offline Component}

Offline, we train a flow matching velocity field $v_\theta(\tau_t,t)$
on $\mathcal{D}$ via conditional flow matching~\cite{lipman2023flow}
with $\tau_t$ linearly interpolating Gaussian noise and a
demonstration. The flow--score identity~\cite{song2021score} yields
$s_\theta\!\approx\!\nabla\log p^{\mathrm{data}}_t$ at inference,
in line with~\cite{flowmp,janner2022}. Amortizing the demonstrations
into a network replaces $O(MH)$ rollouts with one
$O(|\theta|)$ forward pass per denoising step---an offline analogue
of fitted value iteration, paid once at
$O(N\!\cdot\!\text{epochs}\!\cdot\!|\theta|)$. As a proxy until
$v_\theta$ is trained, the present report uses a 5-layer MLP
behavior-cloning policy $\pi_\phi(a\mid o)$ on the same data; this
stand-in lets us measure the offline$\to$online gap ahead of the
full generative model.

\subsection{Online Component}

At inference we evaluate
$\nabla\log p^{\mathrm{model}}_t\!\approx\!
\mathbb{E}_{w(\tau)}[\tau-\tau_t]/\sigma_t^2$ from self-normalized
importance sampling of cost-weighted rollouts
$w(\tau)\propto\exp(-J(\tau)/\beta)$, in the spirit of
MBD~\cite{mbd}, and add it to $s_\theta$ at every step. The offline
prior is blind to the deployed cost and degrades out of
distribution; MBD has high estimator variance over long horizons
and lacks a data prior. We mitigate both by anchoring rollouts at
the flow-matching posterior mean and letting $\lambda(x,t)$ react
to local data density. In the LQ case the correction reduces to
$-P_t x_t$, an LQR-style sanity check we realized this milestone
(\S\ref{sec:initial-results}).

\section{Initial Results}
\label{sec:initial-results}

\subsection{Modeling \& Simulation}

We instantiate the framework on a 7-DoF Franka Panda doing
\texttt{PickPlaceCan} in \texttt{robosuite}/\texttt{MuJoCo}: the
agent picks a cylindrical can and places it in a bin. The dataset's
native action is 7-D operational-space control (3 translation, 3
rotation, 1 gripper) over $H\!\approx\!120$ steps at 20\,Hz, and
demonstrations come from \texttt{Robomimic}'s
\texttt{can}/\texttt{ph} (proficient-human) split. Two pivots from
the Week-3 plan simplified the testbed: we use the dataset's
native OSC actions rather than retargeting to torque control (which
avoids an extra source of distributional drift), and the offline
component is the MLP-BC stand-in (\S IV-A) while $v_\theta$ is
being engineered. As an analytical companion we use a 2D
point-mass LQR with $x_{k+1}=Ax_k+Bu_k$; the discrete algebraic
Riccati equation~\cite{lewis2012optimal} provides closed-form
$K^{\star}$ for validating the online component in isolation.

\subsection{Baseline methods}

We compare five methods on a shared 20 init-states $\times$
3-rollouts protocol (n=60) from the first 20
\texttt{can}/\texttt{ph} demonstrations. \textit{Open-loop demo
replay} replays the recorded action sequence and gives the
empirical \emph{upper bound} (below 1.0 due to documented
\texttt{robosuite}~1.4$\to$1.5 controller-config drift).
\textit{Random} executes $a_t\!\sim\!\mathrm{Uniform}([-1,1]^7)$
as the trivial lower bound. \textit{Behavior cloning} trains
$\pi_\phi$ for 500 epochs (MSE, batch~64) and rolls it out
closed-loop on each demo's native horizon. \textit{Online MBD,
no prior} runs the denoising loop from a zero initial mean for
$K\!=\!10$ iterations of $M\!=\!16$ perturbed rollouts each
($H\!=\!120$, $\lambda\!=\!0.5$, cosine
$\sigma\!:\!0.5\!\to\!0.05$), then executes the planned
trajectory open-loop. \textit{Hybrid prior + MBD} (this
milestone's first instantiation of our method) seeds the same
denoising loop with a deterministic BC rollout instead of a
zero initial mean, and halves $\sigma_{\max}$ to $0.25$ so the BC
structure survives iteration~1. The middle three rows are
stand-ins for the offline-only and online-only ablations of the
eventual flow-matching hybrid; the last row is the actual
product-of-experts pipeline (\S IV) with the MLP-BC stand-in for
$v_\theta$.

\subsection{Initial Results}

\begin{table}[h]
\centering
\small
\caption{Success rate on \texttt{can}/\texttt{ph} over 20 init-states
$\times$ 3 rollouts (n=60) unless noted. LQ companion is a separate
analytical sanity check; bookend methods give the apples-to-apples
ceiling and floor.}
\label{tab:results}
\begin{tabular}{l c l}
\hline
Method & Success rate & Role \\
\hline
Open-loop demo replay (Exp~01)        & 0.90 (n=20)        & upper bound \\
Random actions (Exp~02)               & 0.00               & lower bound \\
MLP behavior cloning (Exp~03)         & 0.067              & offline only \\
Online MBD, no prior (Exp~04)         & 0.00               & online only \\
\textbf{Hybrid BC + MBD (Exp~05)}     & \textbf{0.017}     & \textbf{our method} \\
\hline
LQ analytical companion               & 1.059$\times$ DARE & math sanity \\
\hline
\end{tabular}
\end{table}

Table~\ref{tab:results} summarizes the milestone numbers. Replay's
\textbf{0.90} is the apples-to-apples closed-loop ceiling; random
\textbf{0.00} confirms the task is contact-rich enough that
uniform noise cannot solve it. The MLP-BC stand-in reaches
\textbf{0.067} (4/60), a small but non-trivial step above the
floor that quantifies what an imperfect offline prior buys.
The same denoising loop applied to the 2D LQ companion converges
to within \textbf{1.059} of the DARE optimum (5/5 seeds within
10\%), confirming the score, schedule, and softmax weighting are
implemented correctly; any failure on the can task is therefore
attributable to sim coupling or weighting design rather than to
MBD's algorithmic core.

MBD-no-prior reaches \textbf{0.00} on the full 60-rollout protocol
(78.5\,min wall): every rollout descends the planning cost by
1--12 units, but the executed-reward distribution is dominated by
``cannot reach'' starts (53\,\% with $r\!\le\!+0.5$); only 30\,\%
produce a strong reach ($r\!>\!+3$) and none discover the grasp
in the 7-D$\,\times\,$120-step open-loop search. This is the
predicted failure mode and motivates the hybrid: MBD has no
mechanism to localize the search around demonstration-shaped
behavior.

\textbf{Hybrid BC + MBD} (Exp~05) seeds each denoising trajectory
with a deterministic BC rollout from the same Exp~03 checkpoint
and otherwise uses the Exp~04 loop with a halved
$\sigma_{\max}\!=\!0.25$. On the same 60-rollout protocol
(81.3\,min wall) the headline binary success rate is
\textbf{0.017} (1/60); we treat this as a precision-of-the-
terminal-event artifact rather than as the substantive result,
because Fisher's exact test against Exp~04's 0/60 gives
$p\!=\!1.000$ and the Wilson 95\,\% CI for 1/60 is
$[0.3\%,8.9\%]$, fully overlapping that for 0/60 ($[0\%,6.0\%]$).

The substantive evidence is behavioural, and it is robustly
significant. The ``no reach'' mass
($r_{\mathrm{exec}}\!\le\!+0.5$) collapses from 53\,\% in Exp~04
to \textbf{12\,\%} here ($p\!\approx\!10^{-6}$, Fisher), ``strong
reach'' ($r_{\mathrm{exec}}\!>\!+3$) rises from 30\,\% to
\textbf{53\,\%} ($p\!\approx\!10^{-2}$), and on
\textbf{85\,\% of rollouts} (51/60) the MBD-refined trajectory
executes a strictly higher reward than the BC rollout it was
warm-started from (sign test against 50\,\%:
$p\!\approx\!10^{-8}$). Paired mean executed dense reward rises
from $+4.51$ (BC alone, same protocol) to $\mathbf{+8.25}$
($\mu_{\text{final}}$ after MBD; Wilcoxon signed-rank one-sided
$p\!\approx\!10^{-7}$); the cost descent is broadly positive
(mean $+5.42$, median $+1.16$). These results validate the two
mechanisms predicted by our formulation~\eqref{eq:score}.
(i)~\emph{Reaching is solved by the prior, not by the planner}:
the same MBD search that left 53\,\% of starts motionless from a
zero initial mean leaves only 12\,\% motionless once anchored at
a BC rollout. (ii)~\emph{The planner refines, it doesn't
override}: the 85\,\% win-rate over BC alone shows the
cost-weighted correction is consistently productive rather than
averaging the BC trajectory back to zero. Strict success remains
small because both BC and the hybrid stall on the
millimetre-scale gripper-finger alignment that converts a reach
into a lift; we expect this to close once the flow-matching prior
$v_\theta$ (\S IV-A) replaces the deterministic MLP, since
$v_\theta$ can represent the multimodal demonstration grasp
configurations that the unimodal stand-in cannot.

\section{Next Steps}

In priority order: (i) replace the MLP-BC stand-in with $v_\theta$
and re-run the Exp~05 protocol---the hybrid's bottleneck is the
deterministic, unimodal grasp configuration baked into the MLP,
which a flow-matching prior should fix without changing the
online component; (ii) wire the explicit composite
score~\eqref{eq:score} (Exp~05 currently uses the BC rollout
only as an initial mean rather than as a per-step score term)
and sweep $\lambda\in[0,\infty]$ to map the offline$\,\to\,$online
tradeoff---endpoints should recover the Exp~03 and Exp~04 rows of
Table~\ref{tab:results}; (iii) push $K$ from 10 to 20 and/or add
one MPC-style replan at $t\!=\!60$ to attack the grasp-and-lift
bottleneck that caps Exp~05 at 1/60; (iv) OOD test with novel
goals and $\pm30\%$ mass/friction perturbations; and (v) an
adaptive $\lambda(x,t)$ tied to local data density. We will
revisit \texttt{mujoco-mjx} only if the parameter sweeps in
(ii)--(iii) make the 80-min sequential wall a blocker.

\section{Questions}

We would value feedback on two scoping pivots from the Week-3 plan.
\textbf{(Q1)} Is dropping torque retargeting in favor of the
dataset's native OSC action space acceptable for the
initial-results milestone, given that the apples-to-apples
comparison with all five baselines is preserved? \textbf{(Q2)}
For the hybrid pipeline, is the current ``BC rollout as
$\tau_K$ initial mean'' instantiation a defensible stand-in for
the composite score~\eqref{eq:score} at this milestone, or would
you prefer we hold the headline number until the per-step score
additive (the $\lambda\nabla\log p^{\mathrm{model}}$ term mixed
into every denoising step rather than only at $t\!=\!K$) is in
place?

\bibliographystyle{IEEEtran}
\bibliography{IEEEabrv,references}

\end{document}
